\section{Contracted components and the degree-seven wall formula}\label{p14-sec:virtual}

Let $Z\to C$ be the sextic Galois closure. Put $b=2l$ for the
number of branch values and $h_j=1+7(g-1)+3l+j$.
Fix a nonempty set $I$ of these values. Choose the negative root
line of one inertia transposition $\tau$ in the standard representation.
Over each value in $I$ join the $\tau$-fixed point of $Z$ to the
point of the zero section. Denote the resulting finite map by $a_I^0$
and its signed local degree from Proposition~\ref{p14-prop:signs} by $a_I$.
Its arithmetic genus is $h_{|I|}$. We write $G_{I,k}$ for the locus
obtained from $a_I^0$ by attaching one-pointed contracted curves at
smooth points or inserting two-pointed contracted curves at its nodes,
with total arithmetic genus $k$.

\begin{proposition}\label{p14-prop:virtual-joins}
At a cubic crossing satisfying the kernel and smoothing hypotheses of
Sections~\ref{p14-sec:kernels} and~\ref{p14-sec:joins}, the contribution of
$G_{I,k}$ to the change of the degree-seven local invariant is
\[
       a_I [v^{2k-2}]\bigl(2\sin(v/2)\bigr)^{2h_{|I|}-2}.
\]
A single interval about the crossing suffices to compute the changes
for all $k$ on the same two sides.
\end{proposition}

\subsection{The central locus}
By Section~\ref{p14-sec:kernels}, a central degree-seven map carrying a
projective height has positive normalizations \(Z\) and \(C\), each
with multiplicity one. Its matching height on \(Z\) spans the negative
line of one transposition. Its embedded graph meets the zero section
once over each branch value, with distinct complex tangent lines.

Each connected contracted subcurve at a node or at a smooth point has
one image. Maximal such subcurves have at most two positive branches,
one on each normalization when there are two. The joining values are
therefore recovered from the map with its projective kernel.
For fixed $k$, $G_{I,k}$ is compact. Take stable degree-one maps to
the two normalizations with marks evaluated at the indicated flags.
Glue their marks directly or insert stable two-pointed curves of
positive genus. Summing these proper gluing maps over the finitely
many genus distributions gives the stated locus, including a tail
specializing to a bridge. Distinct joining sets have disjoint images
in the projective incidence. The unique matching height line is
spanned by the chosen root on $Z$ extended by zero on $C$.

\subsection{Displacing the line of sections}
For a stable map $a$ to $C$ write $L_{a,t}=a^*D_t$. Near its
one-dimensional kernel, take a finite-dimensional stabilization and
projectivize its domain over the reals. The zero on the tautological
line represents the local degeneracy class: the derivative in the
projective directions cancels the image directions of the linear map.
In the transverse matrix space, with a generator of the kernel chosen,
the cokernel on the linking sphere is its tangent bundle. The rank off
the central locus is $r=4(k-1-7(g-1))$, the real virtual dimension of
the moduli of maps to $C$. The Euler number of this even-dimensional
link is two. Thus the relative projective class capped with the
fixed-target virtual class times the oriented parameter interval,
and multiplied by two, computes the terminal minus initial Euler
number \cite[\S6.2, Lemma~6.10 and Appendix~B]{bai}.

Choose a closed disk strictly inside a larger disk of regular values.
On the central map the seven sheets over the larger disk are distinct
as pieces of the normalization, including possible contracted
attachments. For nearby maps the preimage of the smaller disk stays
in seven disjoint neighborhoods. Each has degree one by conservation
of degree at an ordinary value avoiding all attaching points. Order
these pieces using an identification of the sextic with the fixed
model over $C$ taking a chosen joining flag to its fixed representative.
There are two such identifications, interchanged by $\tau$. In the
finite covering of all sheet orderings these choices and their
complement are disjoint closed sets over the compact central locus.
Shrinking the neighborhood extends the selection to a stack cover of
degree two, retaining the automorphisms of contracted components.

On one distinguished nonzero sheet choose a real covector density on
the disk with integral $\ell_a(s)\ne0$ on the chosen root $s$.
Integration uses pulled-back area, so constant components give zero.
Normalize by $\ell_a(u)=1$. For independent smooth vector-valued
$(0,1)$-forms on the seven sheets, compactly supported in that disk,
write $j_a$ for their sheetwise pullback. Replace $L_{a,t}u$ by
$L_{a,t}u-\ell_a(u)j_a$. On the normalized slice the equation is
\begin{equation}\label{p14-eq:virtual-affine}
                      L_{a,t}u=j_a,\qquad \ell_a(u)=1 .
\end{equation}
This defines a continuous perturbation in nodal stabilizations.
Indeed introduce $z=\ell_a(u)$, solve $D_t p=j$ locally on each
disk containing a nodal value on each sheet, and replace $u$ by
$u'+za^*p$, using cutoffs equal to one near the nodes. The new forcing
vanishes there, and the scalar row becomes
$z=\ell_a(u')+z\ell_a(a^*p)$. These substitutions are invertible
row and variable changes. On thick parts the remaining forcing converges
smoothly, and on a shrinking disk the area test on bounded sections is
bounded by seven times its area times the supremum of the integrand.
These are precisely the continuous perturbations of the stabilized
inverse in \cite[Appendix~B]{bai} and
\cite[Propositions~B.7.9 and B.10.4]{pardonatlas}.
To obtain a single ambient almost complex structure, split the tangent
bundle of the total space by a complex connection and write
\(D_t=\nabla^{0,1}+A_t\). Over the support disk take disjoint cutoffs
\(\chi_m\) equal to one near the seven normalized sheet graphs.
Write \(F_{t,j}(x,n)=A_t(x)n-\sum_m\chi_m(x,n)j_m(x)\), a fiber-valued
antilinear one-form in the base direction. The block formula
\[
 (X,V)\longmapsto\big(j_CX,\ iV+2iF_{t,j}(x,n)X\big)
\]
squares to minus the identity. With
\(\bar\partial u=(\nabla u+i\nabla u\, j)/2\), its graph equation
near these sheets is \eqref{p14-eq:virtual-affine} before imposing the
scalar row. The construction is independent of the source, including
its contracted components. It extends the homogeneous graph
correspondence of \cite[Lemma~5.9]{bai}. For a compact completion first cut off \(A_t(x)n\) radially,
identically one on a tube containing the normalized central graphs
and their sufficiently small displacements. Use the connection
structure near the divisor at infinity in the projective bundle.
A closed coupling form positive on the projective fibers, plus a
sufficiently large multiple of the base area form, tames this compact
family. All comparisons take place in the smaller tube. In the
bundle tube the area of a closed graph is bounded by its projection
degree times a constant.

On a finite-map chart the sheetwise forms surject onto the cokernel:
a dual solution annihilating them vanishes on an open set of each
positive normalization, and hence everywhere off its nodal flags by
unique continuation. The equation-space dual in the $L^p$ formulation
has no point masses. Sections with prescribed independent values at
the flags then eliminate the dual node-matching rows. Finitely many
forms suffice near the finitely many bare joins. On each smooth
stratum of the finite-map chart the displaced problem is transverse
for a regular choice. A persistent node or a collision of simple
branch values has positive complex codimension, whereas the full
normalized problem with time has index zero. Hence the displaced
zeros are finite, smooth, simply branched covers. Their graphs may
also be taken embedded. A double point over a common value imposes
four real conditions with at most two real parameters for that value.
At a simple ramification point a zero vertical derivative imposes four
conditions. For transversality of these marked problems the same
adjoint test vanishes off the marks. Prescribe arbitrary independent
values and holomorphic first jets there, solving the symbol equation
recursively for the antiholomorphic derivatives of a formal jet.
Testing against these jets eliminates the distributions at the marks.
Immersion excludes the diagonal in the marked-pair test, and finitely
many further forms then suffice by compactness. By conservation of
the local degree their signs in the labelled chart sum to $a_I$.
Here and below all displacements at the finitely many joins are made
using the same ordered disk and orientations.

\subsection{Isolation and contracted components}
Consider a convergent sequence of normalized graph solutions
approaching $G_{I,k}$, with displacements tending to zero. Take stable
limits of the positive normalizations separately, first retaining the
attaching marks. Their nonconstant limits together are the root graph
and the zero section, each once. Every positive component has positive
projection degree, so there are at most two such normalizations.
On forgetting marks, their individual stable limits have no contracted
component. Indeed the theorem of Ekholm and Shende
\cite[Theorem~1.1]{ekholm} would require a vanishing derivative,
equal complex tangent lines at two preimages, or three attaching
preimages of one point. The two central embeddings exclude each
possibility. The theorem applies to these bare maps because the
displaced equation is an ordinary graph equation. For finitely many
smooth perturbation parameters, put the varying structures and the
standard complex structure of an even-dimensional parameter ball
(adjoining one dummy coordinate if needed) in block diagonal form
on the product. Maps in constant slices are holomorphic for this one
structure, with the same attaching tangents.

If there are two positive normalizations, each converges smoothly
to one of the embeddings. Each nearby pair can meet at most once
near each central intersection: projecting the difference of the two
parametrizations to the real four-plane of the limiting tangents has
invertible derivative in the two local source parameters. Away from
these neighborhoods the limiting images are disjoint. Deleting the
one-ended contracted blocks and replacing the two-ended ones by nodes
now gives a finite map satisfying the same equation. Its joining
flags remain distinct. Properness of the marked factors shows that
their limiting subset is $I$. At zero displacement the local isolation
at each bare join forces precisely the central finite map. A regular
displacement has no finite nodal zero there.

With one positive normalization its unmarked limit has direct joins
at a nonempty set $J$. The bare isolation for $J$ again excludes a
sequence of smooth finite zeros at zero displacement. At a regular
displacement these zeros belong to the smooth embedded graphs just
constructed. Each remaining contracted block on such a graph attaches
once. On restoring marks to its limit one can insert rational trees
at smooth points or at existing nodes, and the one-pointed tails
can specialize to these trees. Distinct smooth points over an unjoined
intersection have disjoint fibers under marked stabilization. Hence
no additional joining value results: $J=I$. Conversely, every
degree-one expansion of a displaced graph solves
\eqref{p14-eq:virtual-affine}, and any of its limits at zero displacement
lies in $G_{I,k}$ by the marked description. This proves fixed-genus
isolation and identifies its entire displaced zero set. In particular
a bridge in a limit is counted with its own joining subset.

\subsection{The virtual class at a displaced graph}
Write $\Gamma$ for one of the embedded displaced graphs and $t_0$ for
its parameter. On its smooth projection the undisplaced homogeneous
operator is injective by isolation in the finite-map chart. This
remains true on each degree-one expansion by restriction and matching.
In a chart for maps to $C$ choose stabilizing marks, incidence slices,
horizontal perturbations and obstruction section $s_C$. Use the same
data for maps to the total space, lifting the slices by inverse
image and the horizontal forms by a complex-linear splitting of
projection. The splitting may be chosen affine in height near each
sheet by taking the complex-linear part of a connection lift.
Choose an equivariant supplement $BW$ to the image of $L_{a,t}$ at
the chart center, with representatives on the smooth parts of all
components including contracted ones. On $s_C=0$ the stabilized
vertical equation is
\[
                         L_{a,t}u-Bw=j_a .
\]
Its derivative in $(u,w)$ is invertible. Projection identifies the
regular simultaneous thickening with the horizontal one over the common
family of stabilized domains \cite[Theorem~3.5 and Lemma~5.10]{swam}.
The residual ambient obstruction rows are $s_C$ and $w$. For the
homogeneous projective equation introduce $z=\ell_a(v)$ as before.
Its stabilized solution is $(v,w')=z(u,w)$ and injectivity forces
$z\ne0$ at a projective zero. In the chart $z=1$ the joint rows are
therefore
\[
                       s_C=0,\qquad w=0,\qquad\ell_a(u)-1=0 .
\]
Their Thom product computes both the compactly supported projective
pairing and the constrained ambient class. To compare charts, take
the union of stabilizing marks and direct sums of perturbations.
Adding a vertical supplement pairs its new variable with the same
new obstruction coordinate. Changing horizontal lifts adds rows
vanishing on $s_C=0$, and eliminating a different vertical
supplement adds a multiple of $w$ to the scalar row. These changes
give the same localized Thom product on common refinements
\cite[\S9.2]{pardonatlas}. Both $r$ and the horizontal virtual
dimension are even. The scalar orientation fixed on the smooth
finite-map chart thus applies at every arithmetic genus.

In unparametrized normal coordinates at $\Gamma$ the augmented
linearization with time and the area test is
$(\xi,\dot t)\mapsto(A\xi+B_0\dot t,\lambda\xi)$.
It is invertible by transversality of the normalized finite-map
zero: the horizontal deformation complex of a smooth cover has
torsion cokernel sheaf and no obstruction. The fixed-time operator
$A$ has index zero. If it has a kernel, the augmented isomorphism
makes both kernel and cokernel lines. In an unramified disk where
their generators are nonzero, add a small zeroth-order term linear in
the height minus the height of $\Gamma$. The vertical tangent there
is the normal quotient, so the variation can have nonzero pairing
on these lines. This fixes the graph and makes $A$ invertible,
preserving the augmented isomorphism. In the affine presentation this
changes the homogeneous operator and the forcing by the corresponding
linear and constant terms. Compactness of the degree-one expansions
permits this homotopy inside the isolating chart. Any extra constrained
graphs approaching the retained one would have one nonconstant
normalization: in its retracting tube each positive component projects
in a positive multiple of seven. Ghost censorship at the embedded
multiplicity-one limit gives smooth convergence of the bare
normalization, contrary to uniqueness for the augmented equation.

The nearby ordinary graphs form a regular family $\Gamma_t$.
The same normalization and compactness argument shows that its
fixed-genus locus in a small tube consists of degree-one expansions.
On these maps the area test is the function
$\gamma(t)=\ell(\Gamma_t)-1$, with $\gamma'(t_0)\ne0$ by the
augmented isomorphism. Replacing the scalar row by $\gamma$ preserves
its Thom product with the two other rows since the functions agree
on their zero set. The time slice contributes
$\operatorname{sgn}\gamma'(t_0)$ times the ordinary fixed-time
degree-one class. The normal bundle of $\Gamma$ has degree
$2h_{|I|}-2$, since $c_1(TC)+c_1(N)=0$. Its regular normal
operator pulls back injectively to every degree-one map. Zinger's
comparison identifies the local invariant with the normal Euler class
on the moduli of degree-one maps \cite[Theorem~1.2]{zinger}.
To transport the sign to a complex-linear reference, quotient the
pullback Fredholm complex by the complex on the base curve $\Gamma$.
Restriction to the unique positive normalization splits the pullback
on cokernels, and matching identifies the two kernels throughout any
deformation. The quotient thus has vector-bundle cokernel with its
quotient determinant orientation \cite[Appendix~B]{bai}.
At an invertible end the remaining orientation is the base determinant
sign. At a Dolbeault reference on two non-effective degree
$h_{|I|}-1$ lines the Hodge integral gives the series
$(2\sin(v/2))^{2h_{|I|}-2}$ \cite[Theorem~2]{pandhodge}.
The product of the time and normal signs is the sign of the original
displaced zero, by evaluating in its own smooth genus. Summing these
signs and cancelling the factor two from the Euler link against the
degree two of the labelling proves the asserted coefficient at fixed
$k$. Isotropy on contracted curves is included in the degree-one
virtual class.

\subsection{One interval for exact degree seven}
We record the genus estimate needed to pass to a common interval. For
a compact smooth family of homogeneous normal operators on $C$ there
is a constant $K\ge0$ such that a simple graph projecting in degree
$d$ has normalization genus at most $1+d(g-1)+Kd(d-1)$. Indeed, take
the normalization of the off-diagonal part of the ordered fiber
square of its projection. The difference of the two heights is
generically nonzero on each component. For a fixed unitary connection
and $M$ a uniform norm bound on the zeroth-order term, it satisfies
$|\bar\partial_\nabla s|\le M|s|\,|df|$, up to a fixed metric
constant, for the induced projection $f$. Its isolated zeros have
positive integral orders by the similarity principle. Off these
zeros let $e=s/|s|$ and $\alpha=\operatorname{Im}\langle
e,\nabla e\rangle$. The curvature identity expresses $d\alpha$ as
the curvature evaluated on $e$ plus twice the difference of the squared
norms of the orthogonal $(1,0)$ and $(0,1)$ derivatives of $e$.
The latter $(0,1)$ norm is bounded by $M|df|$. Thus
$d\alpha\ge-f^*\kappa$ for one positive area form $\kappa$ on $C$
uniform over the family. Stokes on the complements of the zeros bounds
their total order by $d(d-1)\int_C\kappa/(2\pi)$.
A ramification point of index $a$ gives $a-1$ distinct off-diagonal
branches over the repeated point, all with zero difference. The genus
bound follows by Riemann--Hurwitz.

Suppose a connected stable map of exact degree $d$ has a nonzero
height. For its positive normalizations let $d_i$ be the degree of
the simple graph and $m_i$ the covering multiplicity, including graphs
of height zero. Thus $\sum_i m_i d_i=d$. Replace each covering by a
connected unramified cover of the simple normalization with the same
multiplicity. Such a cover exists since its base has genus at least
two. The new genus is at most
$1+m_i d_i(g-1)+K m_i d_i(d_i-1)$. Reconnect along a tree of equal
graph values, taken from the original matching after collapsing each
connected contracted block to its image. Lift each selected flag to
its replacement. If flags coincide there, place them on a constant
rational component attached at that point with one node for each of
the coincident flags. This keeps the dual graph a tree and the
resulting map stable. Its degree is $d$, its height is nonzero, and
its genus is at most $dg+Kd(d-1)$. The replacement is used only to
detect existence of a kernel.

Apply this test at $d=7$ on a compact subarc. In each of the finitely
many genera below the bound the central compact loci just described
exhaust all projective zeros. Fixed-genus isolation and compactness
give a kernel-free punctured interval. Intersecting the finitely many
intervals and applying the test proves the assertion in every genus
at exact degree seven. The Euler integrals are then constant on each
punctured side coefficient by coefficient. This finishes the proof of
Proposition~\ref{p14-prop:virtual-joins}. Since the only proper divisor
of seven is one, and the base operator stays invertible, the
degree-one factor has constant exponent. Subtraction of its unchanged
contribution in the connected series gives
$\Delta n_{7,h_j}=\sum_{|I|=j}a_I$ at super-rigid endpoints, with
all other degree-seven exponent changes zero. This proves Theorem~\ref{p14-thm:wall}.
\subsection{The possible changes in elementary multiplicity}\label{p14-sec:chambers}
Write \(c_j=\Delta n_{7,h_j}\). With a common choice of sign as in
Theorem~\ref{p14-thm:wall}, the set of attainable vectors on the given cubic
is the image of the complement of the real hyperplanes
\[
       H_I=\{w\in\RR^b:\, w_I=0\},\qquad
       w_I=\sum_{i\in I}w_i,\qquad \varnothing\ne I\subseteq E ,
\]
under \(w\mapsto ((-1)^j\sum_{|I|=j}\sgn w_I)_{1\leq j\leq b}\).
These hyperplanes form the resonance arrangement \cite{kuehne}. Several
chambers may determine the same vector. The sharp binomial bound used
in Theorem~\ref{p14-thm:torus} follows at once from this description:
\[
       |c_j|\leq\binom bj,\qquad
       \max_{w}\sum_{j=1}^{b}|c_j|=2^b-1 .
\]
Indeed choosing all weights positive gives equality for every \(j\)
simultaneously. Such a choice is realizable by
Proposition~\ref{p14-prop:realize} and the branch variations of
\S\ref{p14-sec:joins}.

\begin{proposition}\label{p14-prop:lattice}
The subgroup generated by the changes on a prescribed cubic is
\[
       2\ZZ^b+\ZZ\left(\binom b1,\ldots,\binom bb\right).
\]
Its index in \(\ZZ^b\) is \(2^{b-1}\).
\end{proposition}
\begin{proof}
Every summand has absolute value one, so the reduction of any
attainable vector modulo two is the displayed binomial vector.
For each nonempty \(I\) choose a point of \(H_I\) outside the finitely
many other hyperplanes. Crossing there changes only \(\sgn w_I\).
The difference of the two attainable vectors is a signed double
coordinate vector in coordinate \(|I|\). These differences give all of
\(2\ZZ^b\); the last binomial coefficient is one, so the reduction
of the generated subgroup has order two.
\end{proof}

There is a more restrictive relation on chambers meeting the full-sum
hyperplane \(H_E\) at a point on no proper-subset hyperplane.
Complementation pairs each proper subset \(I\) with a subset of
opposite sum on \(H_E\). Since \(b=2\ell\), the two cardinalities
have the same parity. On either side of this point it follows that
\[
             c_{b-j}=-c_j\quad(1\leq j<b),\qquad c_\ell=0 .
\]
Only \(c_b\) changes when crossing \(H_E\) there. Consequently neither
of these two chambers attains the unrestricted binomial bound.
For example, at \(b=2\) they give \((0,1)\) and \((0,-1)\). The
same-sign chambers give \((-2,1)\) and \((2,-1)\). Expanding
\(x=v^2-v^4/12+O(v^6)\) gives, for \(g=2\), the connected
genus-thirteen coefficient \(c_2-11c_1/12\). Thus the two chambers
with nonzero \(c_1\) give \(17/6\) and \(-17/6\).

The subgroup from chambers adjacent to \(H_E\) also has a simple
description. In the free coordinates \((c_1,\ldots,c_{\ell-1},c_b)\)
it is
\[
 2\ZZ^\ell+\ZZ\left(\binom b1,\ldots,\binom b{\ell-1},1\right).
\]
Indeed two proper-subset hyperplanes have the same restriction to
\(H_E\) only when their indexing sets agree or are complements:
their indicator vectors must be proportional modulo the constant vector.
Cross one generic restricted hyperplane in \(H_E\), and move both
endpoints slightly to the same side of \(H_E\). This gives twice any
coordinate with \(j<\ell\), up to sign. Crossing \(H_E\) itself gives
twice the last coordinate. Reduction modulo two finishes the argument
also for \(\ell=1\). Here the subgroup describes sums of changes over
several crossings; an individual crossing must satisfy the subset-sign
formula.

\begin{proof}[Proof of Theorem~\ref{thm:invisible}]
Such covers exist by the Riemann existence theorem: send the two peripheral loops to the same transposition and one handle generator to a three-cycle, with its partner trivial. The product relation and transitivity hold. The construction of \S\ref{p14-sec:kernels} realizes the crossing on any prescribed such cover. The standard two-plane of \(S_3\) is the central kernel on its normal closure. The fundamental group of an unramified cover, restricted to the punctured base, contains the normal closure of the peripheral loops. Its image therefore contains the normal closure of a transposition, which is \(S_3\). The two-plane has no invariants. Lemma~\ref{p14-lem:kernelplane} gives an invertible pullback on each unramified cover at the center. There are finitely many covers of the bounded degree, so their determinants remain invertible on a common interval. The smooth cubic calculation of \S\ref{p10-sec:incidence} gives the first change. Theorem~\ref{p14-thm:wall} gives the second from the single subset containing both branch values.
\end{proof}
